\section{Task 5: The demonstration video}
\sectionmark{Task 5: Demonstration video}
\label{sec:t5}

\subsection{Purpose, conditions and criterion}

\begin{plain}{Task 5 at a glance}
\textbf{Purpose.} A video of one minute or less shows the complete system at work: a printed
line is laid below the camera, a key is pressed, the line appears as \LaTeX{}, and after three
lines the typeset PDF opens. The video is handed in through Google Drive, and it is also the
last line of defence of the live demonstration, should the hardware fail on the day.

\textbf{Before you begin.} Task~3 is complete (the pipeline is frozen), and the list of test
lines that were read correctly is available from part~B of Task~4. The recording takes place at
the stand of Task~1. The simplest arrangement is to use the laptop on which Tasks~1 to~4 were
carried out. Any other laptop must have completed the Set-up Guide, except that the two long
steps that generate training images and train the classifier are replaced by \cmd{git pull},
which fetches the team's classifier; and Step~\ref{st:ping} of this document must succeed on
it. A telephone with a camera is needed, and something to hold it still.

\textbf{Complete when.} (1)~The video lasts 60 seconds or less and shows three lines being read
and the resulting PDF. (2)~It is in Google Drive, and its link opens in a browser window that is
not signed in. (3)~The link is entered in the report and has been sent to the team. (4)~A copy
of the file is on the laptop that will be used for the live demonstration.

\textbf{Time.} About an hour and a half, most of it rehearsal.
\end{plain}

\subsection{Choose the three lines}

\step{Select candidates from the results}
The video shows one line of each kind: ordinary text, a formula with superscripts, and a
fraction. If the gate of Task~3 decided that fractions are not demonstrated, the third line is
a second formula. Choose only lines that the frozen pipeline read correctly, that is, lines
marked \texttt{ok} in the final measurement (Step~\ref{st:final}). Table~\ref{tab:cards} lists
suitable lines.

\begin{table}[H]
\centering
\begin{tabular}{C{2.3cm}C{1.3cm}C{3.3cm}C{1.3cm}C{2.0cm}}
\toprule
\hd{Kind} & \hd{Line} & \hd{Content} & \hd{Page} & \hd{Mode} \\
\midrule
Text & T01 & Hello World & 1 & text \\
Text & T03 & ECE 105 Lab & 1 & text \\
Superscripts & T26 & $x^{2}+y^{2}=r^{2}$ & 6 & math \\
Superscripts & T27 & $E=mc^{2}$ & 6 & math \\
Fraction & T54 & $\dfrac{a+b}{2}$ & 11 & math \\[7pt]
Fraction & T47 & $\dfrac{a}{b}+\dfrac{c}{d}$ & 10 & math \\[7pt]
Fraction & T46 & $\dfrac{1}{2}$ & 10 & math \\[5pt]
\bottomrule
\end{tabular}
\caption{Lines suitable for the video and for the live demonstration. The pages are those of
the test lines printed at 40\,pt.}
\label{tab:cards}
\end{table}
\checkpoint{Three lines are chosen, one of each kind, all of them marked \texttt{ok} in the
final measurement. Their pages, printed a second time, are the \emph{cards} of the
demonstration.}

\step{Open the project in PowerShell}
% The opening sequence, written out in full. Input wherever a task may begin in a new window.
If a PowerShell window whose prompt begins with \texttt{(.venv)} is still open, use it and go
on to the next step. Otherwise open a new window (click \ui{Start}, type \code|powershell|,
click \ui{Windows PowerShell}) and give the two commands of the opening sequence in it.
\begin{ps}[Opening sequence, command 1 of 2]
cd C:\ee105\camera-to-latex
\end{ps}
\begin{ps}[Opening sequence, command 2 of 2]
.\.venv\Scripts\Activate.ps1
\end{ps}
\checkpoint{The prompt reads:}
\venvprompt


\step{Confirm each line at the stand}\label{st:confirm}
A line that was read correctly once in a stored photograph may still fail under slightly
different light. Set up the line position, start the application,
\begin{ps}
python -m pipeline.main --new
\end{ps}
and read each of the three lines five times, moving the page a little between readings: press
\key{Space}, compare the reading below the picture with the printed line, and press \key{R}.
Before the text line, press \key{M} until the status line below the picture says \emph{mode
text}; before the formulas, until it says \emph{mode math}. Finish with \key{Q}.
\checkpoint{Each of the three lines was read exactly right at least four times out of five. A
line that fails more often is replaced by another line of Table~\ref{tab:cards} and confirmed
in the same way.}

\subsection{Prepare the scene}

\step{Arrange the screen}
The viewer must be able to read the screen in the video. Three windows matter: the application
window with the picture of the camera, the PowerShell window in which the readings are printed,
and the PDF that opens at the end.
\begin{enumerate}
\item Enlarge everything on the screen: open \menu{Settings}\then\menu{System}\then\menu{Display}
  and, under \ui{Scale}, choose the next larger value, for example 150\,\% in place of
  125\,\%. Note the previous value, so that it can be restored afterwards.
\item Start the application (the command of Step~\ref{st:confirm}). Drag the application window
  to the left half of the screen and the PowerShell window to the right half.
\item Close every window that shows an earlier \file{document.pdf}. Some PDF readers lock the
  file while it is open, and the application could then not write the new one.
\end{enumerate}
\checkpoint{Both windows are visible side by side, and no PDF is open.}

\step{Arrange the camera of the telephone}
Hold the telephone horizontally and fix it, on a tripod or leaning against a stack of books,
about 70\,cm from the desk, so that its picture contains both the stand, with the page below
the board, and the laptop screen. Switch off the flash. Record a test of ten seconds and look
at it.
\checkpoint{In the test recording the page below the board can be seen, the text below the
picture in the application window can be read, and no lamp is reflected in the laptop screen.}

\subsection{Rehearse and record}

\step{Learn the sequence}
Table~\ref{tab:storyboard} gives the sequence of the video with the time that each part may
take. The keys act on the application window, which must be the active window: click on its
picture once before starting.

\begin{table}[H]
\centering
\begin{tabular}{C{1.6cm}L{5.0cm}L{4.3cm}}
\toprule
\hd{Seconds} & \multicolumn{1}{c}{\hd{Action}} & \multicolumn{1}{c}{\hd{Words (optional)}} \\
\midrule
0--6 & Show the stand; point to the board & ``An EE\,105 board photographs one printed line.'' \\
6--20 & Press \key{M} until \emph{mode text}. Lay the first card; wait for the picture; press
  \key{Space}, then \key{A} & ``The laptop cuts the photograph into symbols and names each one.'' \\
20--33 & Press \key{M} for \emph{mode math}. Lay the second card; \key{Space}, then \key{A} &
  ``Positions and sizes give the superscripts.'' \\
33--46 & Lay the third card; \key{Space}, then \key{A} & ``A bar with symbols above and below
  becomes a fraction.'' \\
46--58 & Press \key{P}; the PDF opens & ``The accepted lines are typeset by \LaTeX{} into a PDF.'' \\
\bottomrule
\end{tabular}
\caption{Sequence of the video. After the space bar, wait until the boxes and the reading have
appeared before pressing the key A; after the key A, wait until the live picture has returned.}
\label{tab:storyboard}
\end{table}

What each key does is shown in the PowerShell window: after \key{Space} the reading is printed;
after \key{A} the line \emph{accepted line 1}, followed by the reading; after \key{P} the path
of the PDF, which then opens in the PDF reader of the laptop.

\step{Rehearse with a stopwatch}
Run through the whole sequence three times without recording, with the stopwatch of the
telephone. Begin each run with a fresh session: press \key{Q} in the application window, close
the PDF, and start the application again with the command of Step~\ref{st:confirm}.
\checkpoint{One complete run took 58 seconds or less, and all three readings were correct. If
the runs are too long, omit the spoken words and the first six seconds; if they are still too
long, show two lines.}

\step{Record}
Start a fresh session as before, start the recording on the telephone, and carry out the
sequence. If a reading is wrong, stop and begin again; a video of the demonstration should not
need an excuse. Record until two takes are free of errors, and keep the better one.
\checkpoint{The chosen take lasts 60 seconds or less from its first to its last picture, all
three readings in it are correct, and the PDF at the end shows the three lines. If the take
has a few idle seconds at its beginning or end, trim them with the editing function of the
telephone's photo application; do not cut anything out of the middle.}

\step{Restore the screen}
Set \ui{Scale} back to its previous value (\menu{Settings}\then\menu{System}\then\menu{Display}).

\subsection{Upload and share}

\step{Upload the video to Google Drive}
On the telephone, open the Google Drive application, signed in with the USC account. Tap the
\ui{+} button, then \ui{Upload}, and select the video. Give the upload time to finish on a
wireless network.
\checkpoint{The video appears in \emph{My Drive} and plays when it is tapped.}

\step{Allow others to view it}\label{st:share}
On the laptop, open \url{https://drive.google.com} in the browser, signed in with the same
account. Right-click the video and choose \menu{Share}\then\menu{Share}. Under \ui{General
access}, change \ui{Restricted} to \ui{Anyone with the link}, and leave the role at
\ui{Viewer}. Click \ui{Copy link}, then \ui{Done}.
\checkpoint{The link is in the clipboard.}

\begin{iffail}[If ``Anyone with the link'' is not offered]
The university may restrict sharing. Choose the entry with the name of the university in its
place; every member of the university who is signed in can then view the video. Ask a teammate
to open the link as a check.
\end{iffail}

\step{Verify the link as a stranger would see it}
Open a private window of the browser (\key{Ctrl}+\key{Shift}+\key{N} in Chrome and Edge), paste
the link into the address bar with \key{Ctrl}+\key{V}, and press \key{Enter}.
\checkpoint{The video plays in the private window without any request to sign in. Its length,
shown at the right of the player's time bar, is 1:00 or less.}

\step{Deliver the link and keep a copy}
Paste the link into the place where the course's report template asks for videos, and send it
to the team. Download the video to the laptop that will be used for the live demonstration
(right-click the video in Google Drive, then \menu{Download}), and check that the downloaded
file plays without a network connection.
\checkpoint{The report contains the link, and the file plays on the demonstration laptop with
the wireless network switched off.}

\subsection{If the board is not available}
\label{sec:noboard}

Should the board be unavailable when the video must be recorded, the pipeline can be shown on
the stored photographs of the same three lines. This is a substitute, and the video must say
so: the spoken words, or a caption, state that stored photographs taken earlier with the board
are being read. The three commands below read the photographs of the default lines T01, T26 and
T54; each shows its result in a window until a key is pressed, and the last one builds and
opens the PDF. Each command needs some seconds to start.
\begin{pswide}[Stored photograph of the text line]
python -m pipeline.main --new --mode text --image test_set\images\T01.png --show --accept
\end{pswide}
\begin{ps}[Stored photograph of the line with superscripts]
python -m pipeline.main --mode math --image test_set\images\T26.png --show --accept
\end{ps}
\begin{pswide}[Stored photograph of the fraction, then the PDF]
python -m pipeline.main --mode math --image test_set\images\T54.png --show --accept --pdf
\end{pswide}
\checkpoint{After the third command, the PDF opens with three lines. The same three commands
are the first fallback of the live demonstration (Section~\ref{sec:rehearsal}).}

\subsection{Record}

\begin{table}[H]
\centering
\begin{tabular}{L{6.4cm}C{4.8cm}}
\toprule
\multicolumn{1}{c}{\hd{Quantity}} & \hd{Value} \\
\midrule
Date; persons recording and operating & \\
The three lines shown & \\
Readings correct in the confirmation (of 5, each line) & \\
Length of the video & \hfill s \\
Link checked in a private window & \cbox\ yes\quad\cbox\ no \\
Link entered in the report; copy on the demonstration laptop & \cbox\ yes\quad\cbox\ no \\
\bottomrule
\end{tabular}
\caption{Record of Task 5.}
\label{tab:t5record}
\end{table}
